This paper presents a novel classification framework and the associated deep learning process. Besides the interpretation of the classification task as a regularized optimal transport problem, we demonstrate that this new formalization has some valuable properties about error bounds and structural robustness regarding adversarial attacks. We also propose a systematic approach to ensure the 1-Lipschitz constraint of a neural network. This includes a state-of-the-art regularization algorithm and more precise constant evaluation for convolutional and pooling layers. Even if this regularization process can increase the computation time during learning, it doesn't impact inference. We developed an open source python library based on tensorflow for 1-Lipshitz layers and gradient preserving activation and pooling functions. This makes the approach very easy to implement and to use. \\

The experiment emphasizes the theoretical results and confirms that the classifier has good and predictable robustness to adversarial attacks with an acceptable cost on accuracy. We also show that our classifier forces adversarial attacks to explicitly modify the input. This suggest that our models can use adversarial attacks for explaining a prediction as it is done with counterfactual explanation in logic~\cite{DBLP:journals/corr/abs-1711-00399}.

In conclusion, we believe that this classification framework based on optimal transport is of great interest for critical problems since it provides both empirical and theoretical guarantees. %Future works will focus on the multiclass counterpart of the approach and on its applicability to large and deep networks.
